% ============================================================
% Appendix to the research poster
% "Not Bad, Not Good? Sentiment Models under Negation"
% Compile with: xelatex appendix.tex
% ============================================================
\documentclass[11pt,a4paper]{article}

\usepackage{fontspec}
\usepackage[margin=2.5cm]{geometry}
\usepackage[english]{babel}
\usepackage{microtype}
\usepackage{enumitem}
\usepackage{xcolor}
\usepackage{titlesec}
\usepackage[hidelinks]{hyperref}
\usepackage{fancyhdr}

\setmainfont{Lato}[Ligatures=TeX]
\definecolor{navy}{RGB}{39,60,117}
\titleformat{\section}{\large\bfseries\color{navy}}{\thesection}{0.6em}{}
\titleformat{\subsection}{\normalsize\bfseries}{\thesubsection}{0.6em}{}
\titlespacing*{\section}{0pt}{1.4em}{0.6em}
\setlist[itemize]{leftmargin=1.4em,itemsep=0.35em,topsep=0.4em}

\pagestyle{fancy}\fancyhf{}
\renewcommand{\headrulewidth}{0.4pt}
\lhead{\small\color{navy}Appendix --- Not Bad, Not Good? Sentiment Models under Negation}
\rfoot{\small\thepage}

% Hanging-indent reference entry
\newcommand{\refitem}[1]{\par\noindent\hangindent=1.4em\hangafter=1 #1\par\vspace{0.45em}}

\begin{document}

\begin{center}
{\LARGE\bfseries\color{navy} Appendix}\\[0.6em]
{\large Not Bad, Not Good? Sentiment Models under Negation}\\[0.9em]
{\normalsize Priyadith Hosahalli Nagesh \quad(Student ID: 1857563)}\\[0.3em]
{\normalsize Vismaya Basavaraju Mamatha \quad(Student ID: 1784171)}\\[0.7em]
{\small Trends in Natural Language Processing (3.\ Parallelgruppe), SoSe 2026}\\
{\small Universit\"at Trier \quad\textbullet\quad Examiner: Kai Kugler}\\[0.5em]
{\small Repository: \url{https://github.com/PriyadithGowda/negation-sentiment}}
\end{center}

\vspace{0.8em}

% =====================================================================
\section{References}

All sources consulted for this poster are listed below. Entries are grouped into
academic publications, the datasets and model checkpoints used in the experiments,
and the software libraries and documentation relied upon.

\subsection*{Academic publications}

\refitem{Cohen, J. (1960). A Coefficient of Agreement for Nominal Scales. \emph{Educational and Psychological Measurement}, 20(1), 37--46.}

\refitem{Efron, B. (1979). Bootstrap Methods: Another Look at the Jackknife. \emph{The Annals of Statistics}, 7(1), 1--26.}

\refitem{Hossain, M.\,M., Kovatchev, V., Dutta, P., Kao, T., Wei, E., \& Blanco, E. (2020). An Analysis of Natural Language Inference Benchmarks through the Lens of Negation. In \emph{Proceedings of the 2020 Conference on Empirical Methods in Natural Language Processing (EMNLP)}, 9106--9118. \url{https://aclanthology.org/2020.emnlp-main.732/}}

\refitem{Hutto, C.\,J., \& Gilbert, E. (2014). VADER: A Parsimonious Rule-Based Model for Sentiment Analysis of Social Media Text. In \emph{Proceedings of the Eighth International AAAI Conference on Weblogs and Social Media (ICWSM)}, 8(1), 216--225. \url{https://ojs.aaai.org/index.php/ICWSM/article/view/14550}}

\refitem{Kassner, N., \& Sch\"utze, H. (2020). Negated and Misprimed Probes for Pretrained Language Models: Birds Can Talk, But Cannot Fly. In \emph{Proceedings of the 58th Annual Meeting of the Association for Computational Linguistics (ACL)}, 7811--7818. \url{https://aclanthology.org/2020.acl-main.698/}}

\refitem{Liu, Y., Ott, M., Goyal, N., Du, J., Joshi, M., Chen, D., Levy, O., Lewis, M., Zettlemoyer, L., \& Stoyanov, V. (2019). RoBERTa: A Robustly Optimized BERT Pretraining Approach. \emph{arXiv:1907.11692}. \url{https://arxiv.org/abs/1907.11692}}

\refitem{McNemar, Q. (1947). Note on the Sampling Error of the Difference between Correlated Proportions or Percentages. \emph{Psychometrika}, 12(2), 153--157.}

\refitem{Morris, J., Lifland, E., Yoo, J.\,Y., Grigsby, J., Jin, D., \& Qi, Y. (2020). TextAttack: A Framework for Adversarial Attacks, Data Augmentation, and Adversarial Training in NLP. In \emph{Proceedings of the 2020 Conference on Empirical Methods in Natural Language Processing: System Demonstrations}, 119--126. \url{https://aclanthology.org/2020.emnlp-demos.16/}}

\refitem{Ribeiro, M.\,T., Wu, T., Guestrin, C., \& Singh, S. (2020). Beyond Accuracy: Behavioral Testing of NLP Models with CheckList. In \emph{Proceedings of the 58th Annual Meeting of the Association for Computational Linguistics (ACL)}, 4902--4912. \url{https://aclanthology.org/2020.acl-main.442/}}

\refitem{Socher, R., Perelygin, A., Wu, J., Chuang, J., Manning, C.\,D., Ng, A., \& Potts, C. (2013). Recursive Deep Models for Semantic Compositionality Over a Sentiment Treebank. In \emph{Proceedings of the 2013 Conference on Empirical Methods in Natural Language Processing (EMNLP)}, 1631--1642. \url{https://aclanthology.org/D13-1170/}}

\refitem{Truong, T.\,H., Baldwin, T., Verspoor, K., \& Cohn, T. (2023). Language Models Are Not Naysayers: An Analysis of Language Models on Negation Benchmarks. In \emph{Proceedings of the 12th Joint Conference on Lexical and Computational Semantics (*SEM)}, 101--114. \url{https://aclanthology.org/2023.starsem-1.10/}}

\refitem{Wang, A., Singh, A., Michael, J., Hill, F., Levy, O., \& Bowman, S.\,R. (2018). GLUE: A Multi-Task Benchmark and Analysis Platform for Natural Language Understanding. In \emph{Proceedings of the 2018 EMNLP Workshop BlackboxNLP}, 353--355. \url{https://aclanthology.org/W18-5446/}}

\refitem{Yang, A., Yang, B., Zhang, B., et al.\ (Qwen Team) (2024). Qwen2.5 Technical Report. \emph{arXiv:2412.15115}. \url{https://arxiv.org/abs/2412.15115}}

\subsection*{Datasets and model checkpoints}

\refitem{\textbf{SST-2 (GLUE version)} --- Stanford Sentiment Treebank, binary sentiment task. Dataset card: \url{https://huggingface.co/datasets/stanfordnlp/sst2}. Underlying corpus: Socher et al.\ (2013); benchmark packaging: Wang et al.\ (2018). \emph{Only the validation split was used}, because the RoBERTa checkpoint below was fine-tuned on the training split.}

\refitem{\textbf{textattack/roberta-base-SST-2} --- RoBERTa-base fine-tuned on SST-2, released with the TextAttack framework (Morris et al., 2020). Model card: \url{https://huggingface.co/textattack/roberta-base-SST-2}}

\refitem{\textbf{Qwen/Qwen2.5-3B-Instruct} --- instruction-tuned 3B-parameter language model. Model card: \url{https://huggingface.co/Qwen/Qwen2.5-3B-Instruct}. Technical report: Qwen Team (2024).}

\subsection*{Software and documentation}

\refitem{\textbf{vaderSentiment} (Python implementation of VADER). \url{https://github.com/cjhutto/vaderSentiment}}

\refitem{\textbf{Hugging Face Transformers} --- model loading and inference. \url{https://huggingface.co/docs/transformers}}

\refitem{\textbf{scikit-learn} --- Cohen's $\kappa$ (\texttt{sklearn.metrics.cohen\_kappa\_score}). \url{https://scikit-learn.org}}

\refitem{\textbf{statsmodels} --- exact McNemar test (\texttt{statsmodels.stats.contingency\_tables.mcnemar}). \url{https://www.statsmodels.org}}

\refitem{\textbf{pandas}, \textbf{NumPy}, \textbf{Matplotlib} --- data handling, bootstrap resampling and figures. \url{https://pandas.pydata.org}, \url{https://numpy.org}, \url{https://matplotlib.org}}

\refitem{\textbf{Gemini poster theme} for \LaTeX{} beamerposter, by Anish Athalye (MIT licence), used as the basis for the poster layout. \url{https://github.com/anishathalye/gemini}}

\newpage

% =====================================================================
\section{Use of generative AI tools}

The declaration of academic integrity requires that generative AI tools be identified
by product name, that all tools used be listed in an appendix, and that AI-generated
passages be marked. This section provides that list.

\subsection*{Tools used}

\begin{itemize}
  \item \textbf{Claude} (Anthropic) --- used as a programming and drafting assistant.
  \item \textbf{Gemini} (Google, as integrated in Google Colab) --- used for explanations of Python error messages during setup.
\end{itemize}

\subsection*{What the tools were used for}

\begin{itemize}
  \item \textbf{Code.} The Python in \texttt{scripts/} --- data preparation, the rule-based
        perturbation generator, the model-inference code and the statistical analysis
        (McNemar tests) --- was written with AI assistance. The \LaTeX{} source of the
        poster and of this appendix was likewise produced with AI assistance.
  \item \textbf{Text.} Most prose on the poster and in this appendix was drafted and edited
        with AI assistance. Two passages on the poster are the authors' own work and are
        listed below. Every number reported on the poster was checked against
        \texttt{results/results.json}, and the statistics were recomputed from
        \texttt{results/predictions.csv} by a second, independent script before submission.
\end{itemize}

\subsection*{What the authors did themselves}

\begin{itemize}
  \item Choice of topic, research questions and hypotheses.
  \item Selection of the 44 seed sentences from the SST-2 validation split.
  \item Manual inspection of all 264 drafted rows (44 unmodified seed sentences and
        220 rule-generated perturbations). 34 rows were flagged: 33 were reworded and
        kept, and one was excluded because the antonym substitution had reversed the
        polarity instead of preserving it, leaving 263 items. Three further sentences in
        the same seed group were shortened to match a rewording. The row-by-row record is
        in \texttt{data/verification\_record.csv}.
  \item Independent, blind labelling of an 80-item random sample of the drafts by both
        authors --- before the verification pass and before any discussion --- from which
        Cohen's $\kappa$ was computed.
  \item Interpretation of the results and the conclusions drawn from them.
  \item Rewriting, in their own words, of two passages on the poster: the Discussion \&
        Conclusion block and the caption of Table~2. AI assistance on these two was limited
        to one grammatical correction (``Both H2 and H3 are not upheld'' $\rightarrow$
        ``Neither H2 nor H3 is upheld''), two wording corrections in the caption
        (``Select one representative model'' $\rightarrow$ ``One representative model'';
        ``rules-based'' $\rightarrow$ ``rule-based''), and the deletion of ten words in total
        so that both passages fit the fixed column width of the A1 layout. No wording was
        substituted for the authors' own.
  \item The perturbed sentences in the test suite were first generated by a rule-based
        script over a hand-built antonym lexicon. These drafts are not model-generated
        text: they are the output of explicit substitution rules. (One rule --- the
        choice of contrast clause --- uses Python's \texttt{hash()}, so re-running the
        generator can pick a different clause; the sentences actually used are fixed in
        \texttt{data/test\_suite\_final.csv}.)
\end{itemize}

\vspace{0.4em}
\noindent
\textit{The authors take full responsibility for all content of the poster and this
appendix, including any part produced with AI assistance.}

\vspace{2em}

% =====================================================================
\section{Declaration of Academic Integrity}

\noindent
The signed copy of the official University of Trier declaration of academic integrity
(\emph{Eigenst\"andigkeitserkl\"arung}) is attached on the following page(s).

\vspace{0.6em}
\noindent
Form source: \url{https://www.uni-trier.de/fileadmin/organisation/ABT2/HPA/BA-MA-Arbeit/Eigenstaendigkeitserklaerung_EN.pdf}

\end{document}
